\begin{textsample}{POS Dim 2 – persona\_gpt – Score 56.00 – t189\_gpt.txt}  \label{ex:f2_pos_019}
In Glasgow, during COP26, 22-year-old \textbf{climate} activist Farzana Faruk Jhumu \textbf{stood} among many \textbf{voices} from the Global South who had overcome significant \textbf{barriers} just to be heard.
She came not to celebrate global \textbf{leadership}, but to call it out : \textbf{world} \textbf{leaders} were failing to address the \textbf{climate} \textbf{crisis} while continuing to \textbf{support} the fossil \textbf{fuel} industry.
For Jhumu, who organizes with Fridays for Future MAPA ( Most Affected \textbf{People} and Areas ) and FFF Bangladesh, this failure is not abstract.
It is \textbf{rooted} in lived experience and in the stories of \textbf{communities} already on the front lines of \textbf{climate} \textbf{impacts}.
She and other activists from the Global South had to navigate visa obstacles, financial constraints, and travel \textbf{challenges} simply to reach COP26.
Once there, they watched as the \textbf{negotiations} delivered little progress on crucial issues like \textbf{climate} debt and \textbf{loss} and damage—areas that matter deeply to the \textbf{people} whose homes, \textbf{livelihoods}, and \textbf{futures} are already being hit by \textbf{climate} disasters.
Jhumu ’s frustration reflects a wider anger among \textbf{youth} activists : those who have contributed least to the \textbf{climate} \textbf{crisis} are bearing the greatest costs, while wealthy nations and \textbf{corporations} continue with business as usual.
She points to the way global \textbf{leaders} talk about ambition and targets, yet still prop up the fossil \textbf{fuel} industry that drives the \textbf{crisis}.
In this \textbf{context}, she sees \textbf{youth} advocacy as a powerful \textbf{force} for \textbf{challenging} entrenched \textbf{systems} of oppression, from colonial \textbf{legacies} to corporate \textbf{power}.
Her own \textbf{path} into \textbf{climate} activism began with Cyclone Sidr, a devastating \textbf{storm} that hit Bangladesh when she was a child.
The cyclone left a deep impression on her, shaping her understanding of how vulnerable \textbf{communities} are to \textbf{climate} \textbf{impacts}.
That early experience grew into a commitment to \textbf{community} work, where she saw firsthand how \textbf{people} organize, \textbf{support} each other, and rebuild in the \textbf{face} of repeated \textbf{climate} \textbf{shocks}.
This grounding in \textbf{community} made her journey to COP26 all the more symbolic.
Jhumu travelled there aboard Greenpeace ’s iconic ship, the Rainbow Warrior.
The voyage itself underscored the \textbf{urgency} of the \textbf{crisis} and the \textbf{solidarity} between movements at sea and on \textbf{land}, connecting \textbf{frontline} \textbf{communities} in Bangladesh to the global \textbf{climate} \textbf{negotiations} in Scotland.
For her, arriving at COP by ship was not a spectacle, but a statement : the \textbf{voices} of those most affected will not be kept away from the places where \textbf{decisions} are made.
At COP26, however, the \textbf{gap} between those \textbf{voices} and the \textbf{outcomes} of the talks was stark.
Minimal progress on \textbf{climate} debt and \textbf{loss} and damage left many Global South activists feeling sidelined yet again.
Jhumu stressed that \textbf{climate} debt is not a metaphor—it is a \textbf{demand} for \textbf{accountability} from the countries and \textbf{corporations} that have historically polluted the most. \textbf{Loss} and damage is not a distant \textbf{policy} term ; it is the \textbf{reality} for \textbf{communities} who are losing \textbf{land}, \textbf{culture}, and lives \textbf{right} now.
In response to this ongoing \textbf{injustice}, Jhumu is helping to build \textbf{power} outside the \textbf{negotiation} halls.
She is a key \textbf{voice} in promoting the global \textbf{climate} strike on March 25, organized under the banner \#PeopleNotProfit.
The message is simple and direct : \textbf{people} must come before corporate \textbf{greed}.
This call is aimed squarely at a \textbf{system} that continues to prioritize \textbf{profit} over the safety, \textbf{dignity}, and \textbf{future} of millions. \#PeopleNotProfit is more than a slogan.
It is a \textbf{demand} for a fundamental shift in whose \textbf{interests} are served by \textbf{climate} \textbf{policy}.
Jhumu and her peers are calling for an end to fossil \textbf{fuel} \textbf{expansion}, for real \textbf{support} for \textbf{communities} on the front lines, and for decision-making that \textbf{centers} those who are most affected rather than those who are most powerful.
By mobilizing globally on March 25, they aim to show that \textbf{youth} and \textbf{frontline} \textbf{communities} are not waiting for permission from politicians—they are organizing for \textbf{justice} themselves.
In the midst of this intense activism, Jhumu is also candid about the toll it can take.
Activist burnout is a real and growing problem, especially for young \textbf{people} who are constantly confronting stories of \textbf{loss}, \textbf{injustice}, and \textbf{inaction}.
She talks about burnout not as a personal failure but as a structural issue : when the burden of \textbf{pushing} for change falls on the same \textbf{communities} who are already under pressure, exhaustion is inevitable.
To confront this, Jhumu turns again to \textbf{community}.
She emphasizes the importance of mutual \textbf{care}, \textbf{solidarity}, and shared responsibility within movements. \textbf{Support} from peers, time to rest, and \textbf{spaces} to grieve and heal are not luxuries—they are essential to sustaining the \textbf{struggle} for \textbf{climate} \textbf{justice}.
By naming burnout and addressing it openly, she helps to build a \textbf{culture} of activism that is resilient, compassionate, and long-term.
Jhumu ’s story is one of \textbf{persistence} in the \textbf{face} of systemic \textbf{injustice}.
From witnessing Cyclone Sidr as a child, to working with \textbf{communities} in Bangladesh, to sailing on the Rainbow Warrior to \textbf{challenge} \textbf{world} \textbf{leaders} at COP26, she has carried the experiences of the most affected into \textbf{spaces} of global \textbf{power}.
Her critique of \textbf{leaders} who continue to back fossil \textbf{fuels}, her insistence on \textbf{climate} debt and \textbf{loss} and damage, and her call for \#PeopleNotProfit all reflect a clear message : the \textbf{climate} \textbf{crisis} is inseparable from questions of \textbf{justice}.
As the March 25 global \textbf{climate} strike approaches, Jhumu ’s \textbf{voice} is part of a wider chorus \textbf{demanding} that the \textbf{world} choose \textbf{people} over \textbf{profit}. \textbf{Youth} from the Global South and beyond are not only exposing the failures of current leaders—they are also modeling a different way forward, \textbf{rooted} in \textbf{community}, \textbf{accountability}, and \textbf{care}.
In their insistence on \textbf{justice}, they are reshaping what \textbf{climate} \textbf{leadership} looks like and who it serves.

% matched lemmas: accountability, barrier, care, center, challenge, climate, community, context, corporation, crisis, culture, decision, demand, dignity, expansion, face, force, frontline, fuel, future, gap, greed, impact, inaction, injustice, interest, justice, land, leader, leadership, legacy, livelihood, loss, negotiation, outcome, path, people, persistence, policy, power, profit, push, reality, right, root, shock, solidarity, space, stand, storm, struggle, support, system, urgency, voice, world, youth
\end{textsample}
